\section*{अध्याय \ref{ch:g11:magnetic-fields} --- चुंबकत्व और चुंबकीय क्षेत्र}

\begin{solution}{exo:g11:magnetic-fields:1}
हर टुकड़ा उत्तरी और दक्षिणी ध्रुव वाला पूरा \omterm{def:g11:magnetic-fields:magnet}{चुंबक} है: हर कटे फलक पर एक
नया ध्रुव उग आता है। ताज़े फलक एक N और एक S हैं, इसलिए वे आकर्षित करते हैं ---
टुकड़े फिर से जुड़ना चाहते हैं। कोई तरीक़ा ध्रुव को अलग नहीं कर सकता: हर कट
एक युग्म ही बनाता है।
\end{solution}

\begin{solution}{exo:g11:magnetic-fields:2}
(क) ग़लत: हर बिंदु पर क्षेत्र की एक ही दिशा होती है। (ख) सही।
(ग) सही: हर रेखा \omterm{def:g11:magnetic-fields:magnet}{चुंबक} के शरीर से होकर बंद हो जाती है। (घ) ग़लत: सुई
रेखा के \emph{अनुदिश} ठहरती है, उसकी स्पर्शी बनकर।
\end{solution}

\begin{solution}{exo:g11:magnetic-fields:3}
$B = \num{2e-7} \times 10 / 0.020 = \qty{1.0e-4}{T}$ --- पृथ्वी के क्षेत्र का दुगुना।
\end{solution}

\begin{solution}{exo:g11:magnetic-fields:4}
$n = 800/0.40 = \qty{2000}{m^{-1}}$; $B = \num{1.26e-6} \times 2000 \times 1.5 \approx \qty{3.8}{mT}$। $B \propto I$: धारा दुगुनी करने पर \qty{7.5}{mT} मिलता है।
\end{solution}

\begin{solution}{exo:g11:magnetic-fields:5}
$F = ILB = 5.0 \times 0.10 \times 0.20 = \qty{0.10}{N}$ --- यानी $m = F/g = 0.10/9.81 \approx \qty{10}{g}$ का भार।
\end{solution}

\begin{solution}{exo:g11:magnetic-fields:6}
\emph{1.} उसका रास्ता सच्चे उत्तर से $\ang{8}$ पश्चिम है: $2000 \times \sin\ang{8} \approx \qty{2.8e2}{m}$ पश्चिम की ओर।

\emph{2.} सुई के उत्तर से $\ang{8}$ \emph{पूरब} हटकर चलिए।
\end{solution}

\begin{solution}{exo:g11:magnetic-fields:7}
$d = \mu_0 I / (2\pi B) = \num{2e-7} \times 20 / \num{5.0e-5}
= \qty{8.0}{cm}$। ऐसी केबल के एक डेसीमीटर भीतर कुतुबनुमा ग्रह नहीं, तार पढ़ता है ---
उसे जीवित चालकों से दूर रखिए।
\end{solution}

\begin{solution}{exo:g11:magnetic-fields:8}
$n = B/(\mu_0 I) = 0.010 / (\num{1.26e-6} \times 4.0)
\approx \qty{2.0e3}{m^{-1}}$; \qty{25}{cm} पर $N = 1989 \times 0.25 \approx 500$ फेरे।
\end{solution}

\begin{solution}{exo:g11:magnetic-fields:9}
$F = 10 \times 0.12 \times 0.50 = \qty{0.60}{N}$; $a = F/m = 0.60/0.020 = \qty{30}{m/s^2}$ --- लगभग $3g$: छड़ सचमुच उछल पड़ती है।
\end{solution}

\begin{solution}{exo:g11:magnetic-fields:10}
संतुलन: $ILB = mg$, इसलिए $I = \num{5.0e-3} \times 9.81 / (0.20 \times 0.10) \approx
\qty{2.5}{A}$। बल ऊपर की ओर होना चाहिए; $\vect B$ क्षैतिज
और तार के लंबवत होने पर चपटा दायाँ हाथ (हथेली ऊपर) धारा की ज़रूरी दिशा तय कर
देता है।
\end{solution}

\begin{solution}{exo:g11:magnetic-fields:11}
चपटा दायाँ हाथ: (क) पन्ने के ऊपर की ओर; (ख) दाईं ओर; (ग) कोई बल
नहीं --- तार क्षेत्र के समांतर है।
\end{solution}

\begin{solution}{exo:g11:magnetic-fields:12}
\emph{1.} $B = \num{2e-7} \times 5.0 / 0.030 \approx
\qty{33}{\micro T}$; रेखाएँ तार के चारों ओर घूमती हैं, इसलिए ठीक उसके नीचे
क्षेत्र क्षैतिज और तार के लंबवत है: पूरब--पश्चिम।

\emph{2.} सुई सदिश योग के पीछे चलती है: $\tan\theta = 33/20 =
1.67$, इसलिए उत्तर से $\theta \approx \ang{59}$ हटकर।
\end{solution}

\begin{solution}{exo:g11:magnetic-fields:13}
\emph{1.} $I = U/R = 12/2.0 = \qty{6.0}{A}$; $n = 400/0.20 = \qty{2000}{m^{-1}}$; $B_0 = \num{1.26e-6} \times 2000 \times 6.0 \approx \qty{15}{mT}$; \omterm{def:g10:refraction:fiber}{क्रोड} सहित $100 B_0 \approx \qty{1.5}{T}$।

\emph{2.} $P = UI = 12 \times 6.0 = \qty{72}{W}$।

\emph{3.} धारा नहीं तो क्षेत्र नहीं --- और नरम लोहा चुंबकित बना नहीं
रहता: जान-बूझकर बनाई गई इसी बात से बोझ उसी क्षण छूट जाता है।
\end{solution}

\begin{solution}{exo:g11:magnetic-fields:14}
\emph{1.} हर भुजा पर $F = NILB = 50 \times 2.0 \times 0.050 \times 0.30 =
\qty{1.5}{N}$। दोनों भुजाओं में धारा उलटी दिशाओं में बहती है,
इसलिए बल विपरीत हैं: एक बलयुग्म, जो फंदे को घुमा देता है।

\emph{2.} $\vect B$ के समांतर होने पर तार पर कोई \omterm{prop:g11:magnetic-fields:laplace}{लाप्लास बल} नहीं लगता।

\emph{3.} आधे चक्कर बाद बलयुग्म उलट जाता और घुमाव को पलट देता; इसलिए
\omterm{rem:g11:magnetic-fields:motor}{दिक्परिवर्तक} हर आधे चक्कर पर धारा उलट देता है और फंदा एक ही ओर घूमता
रहता है।
\end{solution}

\begin{solution}{exo:g11:magnetic-fields:15}
\emph{1.} $d = \num{2e-7} \times 100 / 1 = \qty{2.0e-5}{m} =
\qty{20}{\micro m}$ --- यानी मिलीमीटर भर मोटे ताँबे के \emph{भीतर} गहरे। अकेले
तार के पास पहुँच में आने वाला कोई बिंदु \qty{1}{T} तक नहीं पहुँचता: प्रबल
क्षेत्र हज़ारों फेरे चढ़ाकर (यानी \omterm{def:g11:magnetic-fields:solenoid}{परिनालिका} से) बनाए जाते
हैं, और साथ में लोहा या अतिचालक।

\emph{2.} $\num{e8} / \num{5e-5} = \num{2e12}$: दस की लगभग बारह घातें।
\end{solution}

\begin{solution}{pb:g11:magnetic-fields:1}
\textbf{1.} सुई स्वयं एक छोटा \omterm{def:g11:magnetic-fields:magnet}{चुंबक} है; पृथ्वी का क्षेत्र उसे
पंक्ति में बिठा देता है, उत्तरी सिरा $\vect B$ के क्षैतिज घटक के अनुदिश, यानी
चुंबकीय उत्तर की ओर।

\textbf{2.} $B = \sqrt{20^2 + 44^2} = \sqrt{2336} \approx
\qty{48}{\micro T}$, क्षैतिज से $\tan^{-1}(44/20) \approx \ang{66}$ नीचे झुका हुआ।

\textbf{3.} रास्ता सच्चे उत्तर से $\ang{6}$ पश्चिम जाता है: $30 \times \sin\ang{6} \approx \qty{3.1}{km}$ पश्चिम।

\textbf{4.} त्रुटियाँ जुड़ती हैं: $\ang{10}$, इसलिए $30 \times \sin\ang{10} \approx
\qty{5.2}{km}$ --- दस डिग्री के लिए
पाँच किलोमीटर।

\textbf{5.} क्षैतिज घटक: ध्रुव के पास क्षेत्र लगभग \omterm{def:g10:universal-gravitation:weight}{ऊर्ध्वाधर} हो जाता
है, क्षैतिज हिस्सा सिकुड़कर शून्य की ओर चला जाता है, और सुई के पास पंक्ति
में बैठने को कुछ बचता ही नहीं।

\textbf{6.} $n = 600/0.30 = \qty{2000}{m^{-1}}$।

\textbf{7.} $I = U/R = 12/1.5 = \qty{8.0}{A}$।

\textbf{8.} $B_0 = \mu_0 n I = \num{1.26e-6} \times 2000 \times 8.0
\approx \qty{20}{mT}$।

\textbf{9.} $50 \times \qty{20}{mT} = \qty{1.0}{T}$ --- नियोडिमियम के ध्रुव पर मिलते क्षेत्र का
दुगुना, और वह भी कहीं बड़े फलक पर।

\textbf{10.} $P = UI = 12 \times 8.0 = \qty{96}{W}$; \qty{10}{\minute} में: $E = 96 \times 600 \approx \qty{58}{kJ}$।

\textbf{11.} क्योंकि वह छोड़ भी देता है: स्विच खोलिए, क्षेत्र मर जाता
है, कार ठीक वहीं गिरती है जहाँ चाहिए। स्थायी \omterm{def:g11:magnetic-fields:magnet}{चुंबक} तो हमेशा जकड़े
रहता है।

\textbf{12.} $F = NILB = 100 \times 3.0 \times 0.040 \times 0.25 =
\qty{3.0}{N}$।

\textbf{13.} एक बलयुग्म: अक्ष के आमने-सामने की भुजाओं पर लगते दोनों विपरीत
बल फंदे को एक ही ओर घुमाते हैं। बराबर पर \emph{समांतर} बल
फंदे को बस बग़ल में सरका देते, घुमाते नहीं।

\textbf{14.} कोई नहीं: क्षेत्र के समांतर तार पर कोई \omterm{prop:g11:magnetic-fields:laplace}{लाप्लास
बल} नहीं लगता।

\textbf{15.} वह हर आधे चक्कर पर फंदे में धारा उलट देता है, ठीक उसी क्षण जब
बलयुग्म का चिह्न बदलने वाला होता है --- इसलिए बलाघूर्ण हमेशा एक ही घुमाव
चलाता रहता है।

\textbf{16.} $N$ दुगुना कीजिए, या $I$ दुगुना, या $B$ दुगुना:
बल का युग्म $F = NILB$ इन तीनों में से हर एक के समानुपाती है।

\textbf{17.} $1.5 / \num{5.0e-5} = \num{3.0e4}$: तीस हज़ार पृथ्वियाँ।

\textbf{18.} $n = B/(\mu_0 I) = 1.5 / (\num{1.26e-6} \times 500)
\approx \qty{2.4e3}{m^{-1}}$ --- और उन फेरों में से हर एक \qty{500}{A} ले जाता है।

\textbf{19.} $P = RI^2 = 0.20 \times 500^2 = \qty{50}{kW}$ --- पूरे मोहल्ले भर की तापन-शक्ति। असली लपेट अतिचालक है:
शून्य प्रतिरोध, शून्य क्षय, बशर्ते उसे परम शून्य से कुछ ही डिग्री
ऊपर रखा जाए।

\textbf{20.} कुतुबनुमा $\num{5e-5}$, क्रेन $1$, MRI $1.5$: जहाज़ दोनों
मशीनों से पाँच पायदान नीचे बैठता है, और वे दोनों एक ही दहाई साझा करती हैं ---
अंतरतारकीय $\num{e-10}$ और न्यूट्रॉन तारे के $\num{e8}$ के बीच। बदला कुछ नहीं, बस
\omterm{def:g11:circuits-and-power:current}{ऐंपियर} (और वे फेरे, लोहा तथा अतिचालक जो उन्हें कई गुना कर देते हैं):
एक क्षेत्र, तीन काम।
\end{solution}
